\documentclass[12pt]{article}
\usepackage[a4paper,margin=1in]{geometry}
\usepackage[utf8]{inputenc}
\usepackage[T1]{fontenc}
\usepackage{amsmath,amssymb}
\usepackage{textcomp}
\usepackage{parskip}
\begin{document}
\begin{center}
{\Large \textbf{2019 STEP I}}
\end{center}

\subsection*{Page 1}
\begin{flushleft}
\bigskip\textbf{Section A: Pure Mathematics}\par
\bigskip\textbf{Question 1.} A straight line passes through the fixed point (1 , k) and has gradient tan θ, where k > 0\par
1 -\par
\[and 0 < θ < \pi. Find, in terms of θ and k, the coordinates of the points X and Y where\]
the line meets the x-axis and the y-axis respectively.\par
\textbf{(i)} Find an expression for the area A of triangle OXY in terms of k and θ. (The point O\par
is the origin.)\par
You are given that, as θ varies, A has a minimum value. Find an expression in terms\par
of k for this minimum value.\par
\textbf{(ii)} Show that the length L of the perimeter of triangle OXY is given by\par
\[L = 1 + tan θ + sec θ + k(1 + cot θ + cosec θ).\]
You are given that, as θ varies, L has a minimum value. Show that this minimum\par
\[value occurs when θ = α where\]
1 cos α\par
\[- = k.\]
1 sin α\par
-\par
Find and simplify an expression for the minimum value of L in terms of α.\par
2 3\par
\bigskip\textbf{Question 2.} The curve C is given parametrically by the equations x = 3t , y = 2t . Show that the\par
2 3 3\par
\[equation of the tangent to C at the point (3p , 2p ) is y = px p .\]
-\par
2 3\par
Find the point of intersection of the tangents to C at the distinct points (3p , 2p ) and\par
2 3\par
(3q , 2q ). Hence show that, if these two tangents are perpendicular, their point of inter-\par
\[section is (u + 1 , u), where u = p + q.\]
-\par
\[The curve C is given parametrically by the equations x = u + 1, y = u. Find the\]
-\par
coordinates of the points that lie on both C and C.\par
\[\leq\]
\[Sketch C and C on the same axes. \leq\]
\[\leq\]
\bigskip\textbf{Question 3.} By first multiplying the numerator and the denominator of the integrand by (1 sin x),\par
-\par
evaluate\par
\[1 \pi\]
4 1\par
dx .\par
\[\geq0 1 + sin x\]
Evaluate also:\par
\[1 \pi 1 \pi\]
3 1 3 1\par
dx and dx .\par
\[\geq 1 \pi 1 + sec x \geq0 (1 + sin x)2\]
\end{flushleft}
\subsection*{Page 2}
\begin{flushleft}
\bigskip\textbf{Question 4.} (i) Find integers m and n such that 3 + 2\sqrt2 = m + n\sqrt2.\par
(\par
4 2\par
\textbf{(ii)} Let f(x) = x 10x + 12x 2. Given that the equation f(x) = 0 has four real roots,\par
- -\par
explain why f(x) can be written in the form\par
2 2\par
\[f(x) = (x + sx + p)(x sx + q)\]
-\par
for some real constants s, p and q, and find three equations for s, p and q.\par
Show that\par
2 2 2 2\par
\[s (s 10) + 8s 144 = 0\]
- -\par
and find the three possible values of s .\par
\[Use the smallest of these values of s to solve completely the equation f(x) = 0 ,\]
simplifying your answers as far as you can.\par
\bigskip\textbf{Question 5.} (i) The four points P , Q, R and S are the vertices of a plane quadrilateral.\par
\[What is the geometrical shape of P QRS if -P-\toQ = -S\toR ?\]
\[What is the geometrical shape of P QRS if -P-\toQ = -S\toR and -P-\toQ = -P\toS ?\]
| | | |\par
\textbf{(ii)} A cube with edges of unit length has opposite vertices at (0 , 0 , 0) and (1 , 1 , 1). The\par
points\par
P (p , 0 , 0), Q (1 , q , 0), R (r , 1 , 1) and S (0 , s , 1)\par
lie on edges of the cube. Given that the four points lie in the same plane, show that\par
\[rq = (1 s)(1 p).\]
- -\par
(a) Show that -P-\toQ = S-\toR if and only if the centroid of the quadrilateral P QRS is at\par
the centre of the cube.\par
Note: the centroid of the quadrilateral P QRS is the point with position vector\par
(-O-\toP + -O-\toQ + -O-\toR + -O\toS),\par
where O is the origin.\par
\[(b) Given that -P-\toQ = -S\toR and -P-\toQ = -P\toS , express q, r and s in terms of p.\]
| | | |\par
Show that\par
4p 1\par
\[cos P QR = - .\]
\[5 4p + 8p2\]
-\par
Write down the values of p, q, r and s if P QRS is a square, and show that the\par
length of each side of this square is greater than .\par
\end{flushleft}
\subsection*{Page 3}
\begin{flushleft}
\bigskip\textbf{Question 6.} In both parts of this question, x is real and 0 < θ < \pi.\par
\textbf{(i)} By completing the square, find in terms of θ the minimum value as x varies of\par
9x 12x cos θ + 4.\par
-\par
2 4\par
Find also the maximum value as x varies of 12x sin θ 9x .\par
-\par
Hence determine the values of x and θ that satisfy the equation\par
4 2\par
\[9x + (9 12 sin θ)x 12x cos θ + 4 = 0.\]
- -\par
\textbf{(ii)} Sketch the curve\par
x\par
\[y = ,\]
x θ\par
-\par
2 2\par
x x\par
\[where θ is a constant. Deduce that either \leq 0 or \geq 4θ.\]
x θ x θ\par
- -\par
By considering the numerator and denominator separately, or otherwise, show that\par
2 2\par
sin θ cos x\par
\[\leq 1.\]
\[1 + cos2 θ sin 2 x\]
Hence determine the values of x and θ that satisfy the equation\par
2 2 2\par
x sin θ cos x\par
\[= .\]
\[4θ(x θ) 1 + cos2 θ sin 2 x\]
-\par
\end{flushleft}
\subsection*{Page 4}
\begin{flushleft}
\bigskip\textbf{Question 7.} Consider the following steps in a proof that \sqrt2 + \sqrt3 is irrational.\par
1. If an integer a is not divisible by 3, then a = 3k 1, for some integer k. In both cases,\par
2 ±\par
a is one more than a multiple of 3.\par
a\par
2. Suppose that \sqrt2 + \sqrt3 is rational, and equal to , where a and b are positive integers\par
b\par
with no common factor greater than one.\par
4 4 2 2\par
\[3. Then a + b = 10a b .\]
4. So if a is divisible by 3, then b is divisible by 3.\par
\[5. Hence \sqrt\{2\} + \sqrt\{3\} is irrational.\]
\textbf{(i)} Show clearly that steps 1, 3 and 4 are all valid and that the conclusion 5 follows from\par
the previous steps of the argument.\par
\textbf{(ii)} Prove, by means of a similar method but using divisibility by 5 instead of 3, that\par
\[\sqrt\{6\} + \sqrt\{7\} is irrational.\]
Why can divisibility by 3 not be used in this case?\par
\bigskip\textbf{Question 8.} The function f is defined, for x > 1 , by\par
x\par
t 1\par
\[f(x) = - dt.\]
\[\geq1 ) t + 1\]
Do not attempt to evaluate this integral.\par
\textbf{(i)} Show that, for x > 2 ,\par
x\par
u 2\par
\[- du = 2 f( x).\]
\[\geq2 ) u + 2 2\]
\textbf{(ii)} Evaluate in terms of f, for x > 0,\par
x\par
u\par
du.\par
\[\geq0 ) u + 4\]
\textbf{(iii)} Evaluate in terms of f, for x > 5,\par
x\par
u 5\par
- du.\par
\[\geq5 ) u + 1\]
\textbf{(iv)} Evaluate in terms of f\par
2 2\par
u\par
du.\par
\[\geq1 \sqrt\{u^2\} + 4\]
\end{flushleft}
\subsection*{Page 5}
\begin{flushleft}
\bigskip\textbf{Section B: Applied Mathematics}\par
\bigskip\textbf{Question 9.} A box has the shape of a uniform solid cuboid of height h and with a square base of side b,\par
where h > b. It rests on rough horizontal ground. A light ladder has its foot on the ground\par
and rests against one of the upper horizontal edges of the box, making an acute angle of α\par
with the ground, where h = b tan α. The weight of the box is W . There is no friction at the\par
contact between ladder and box.\par
A painter of weight kW climbs the ladder slowly. Neither the base of the ladder nor the\par
box slips, but the box starts to topple when the painter reaches height λh above the ground,\par
\[where λ < 1 .\]
Show that:\par
\textbf{(i)} R = kλW cos α , where R is the magnitude of the force exerted by the box on the\par
ladder;\par
\textbf{(ii)} 2kλ cos 2α + 1 = 0 ;\par
sin 2α\par
\textbf{(iii)} μ \geq , where μ is the coefficient of friction between the box and the ground.\par
1 3 cos 2α\par
-\par
\bigskip\textbf{Question 10.} In this question, the x-axis is horizontal and the positive y-axis is vertically upwards.\par
A particle is projected from the origin with speed u at an angle α to the vertical. The particle\par
\textasciicircum{}\circ\par
\[passes through the fixed point (h tan β , h), where 0 < β < 90 and h > 0.\]
\textbf{(i)} Show that\par
2 2\par
\[c ck cot β + 1 + k cot β = 0, ( )\]
- ∗\par
2u\par
\[where c = cot α and k = .\]
gh\par
You are given that there are two distinct values of α that satisfy equation ( ). Let α1\par
∗\par
and α2 be these values.\par
(a) Show that\par
\[cot α1 + cot α2 = k cot β.\]
Show also that\par
\[α1 + α2 = β.\]
(b) Show that\par
\[k > 2(1 + sec β).\]
\textbf{(ii)} By considering the greatest height attained by the particle, show that k \geq 4 sec 2 α.\par
\end{flushleft}
\subsection*{Page 6}
\begin{flushleft}
11 (i) Two people adopt the following procedure for deciding where to go for a cup of tea:\par
either to a hotel or to a tea shop. Each person has a coin which has a probability p of\par
\[showing heads and q of showing tails (where p+q = 1). In each round of the procedure,\]
both people toss their coins once. If both coins show heads, then both people go to the\par
hotel; if both coins show tails, then both people go to the tea shop; otherwise, they\par
continue to the next round. This process is repeated until a decision is made.\par
Show that the probability that they make a decision on the nth round is\par
2 2 n-1\par
(q + p )(2qp) .\par
Show also that the probability that they make a decision on or before the nth round\par
is at least\par
-\par
2n\par
whatever the value of p.\par
\textbf{(ii)} Three people adopt the following procedure for deciding where to go for a cup of tea:\par
either to a hotel or to a tea shop. Each person has a coin which has a probability p\par
\[of showing heads and q of showing tails (where p + q = 1). In the first round of the\]
procedure, all three people toss their coins once. If all three coins show heads, then all\par
three people go to the hotel; if all three coins show tails, then all three people go to\par
the tea shop; otherwise, they continue to the next round.\par
In the next round the two people whose coins showed the same face toss again, but\par
the third person just turns over his or her coin. If all three coins show heads, then all\par
three people go to the hotel; if all three coins show tails, then all three people go to\par
the tea shop; otherwise, they go to the third round.\par
Show that the probability that they make a decision on or before the second round is\par
at least , whatever the value of p.\par
\end{flushleft}
\end{document}
